\documentclass{nfr}

\bibliographystyle{apalike}

\title{Supporting Statement -- Postdoctoral Researcher in the Political Economy of Net Zero}
\author{Zhongli Wang}

\begin{document}

\maketitle

My background is in economics, computational social science, and policy-oriented modeling. My doctoral research examines how firms, consumers, and policy-makers interact in complex economic environments. While the empirical focus of my dissertation is on digital and ICT markets, its broader analytical questions are closely connected to this role: how governance structures shape firm behavior, how business models influence market outcomes, how heterogeneous agents respond to policy constraints, and how institutional design affects welfare, legitimacy, and long-term transformation.

\section*{Relevant Research Background and Analytical Skills}

I am close to completion of my PhD. My dissertation, \textit{Strategic Behavior and Economic Dynamics in ICT Markets: An Agent-based Framework}, develops computational and empirical approaches to studying market dynamics under regulation, consumer heterogeneity, network effects, and strategic firm investment. The dissertation consists of three single-authored papers on data governance, capacity investment, and parameter estimation in agent-based economic models.

The relevance of this training to the political economy of net zero lies in its focus on complex systems, institutional constraints, and strategic interaction. Net-zero transitions are not only technical processes of replacing high-carbon technologies with low-carbon ones. They also involve conflicts over ownership, governance, investment incentives, affordability, public legitimacy, and the distribution of costs and benefits. My PhD has trained me to analyze multi-actor environments in which outcomes emerge from the interaction between firms, consumers, policy rules, and institutional arrangements.

My research demonstrates strong analytical skills across theory, simulation, and empirical validation. In my job market paper, I develop a computational model of firm capacity investment under network effects and congestion. The paper studies how firms make strategic investment and price decisions when consumer demand is shaped by peer influence, service quality, and market feedback. This experience is transferable to net-zero sectors such as energy, housing, transport, and food, where investment decisions are shaped by market structure, consumer behavior, infrastructure constraints, and regulatory incentives.

A key strength I would bring to the project is the ability to analyze how micro-level decisions generate system-level outcomes. In net-zero transitions, it is necessary to examine how households, firms, investors, and local authorities respond to incentives and constraints. My experience with agent-based modeling would allow me to contribute to conceptual and methodological work on how different governance arrangements, ownership structures, and business models influence transition pathways.

\section*{Fit with the Political Economy of Net Zero Project}

The role's focus on ownership structures, governance systems, business models, and value logics strongly matches my interest in how institutional arrangements shape economic outcomes. My paper on data governance examines how regulatory intensity affects firm competition, consumer welfare, and regional policy coordination. The central issue is how governance rules alter firm incentives and social outcomes. This logic is highly relevant to net-zero transitions, where regulation, public ownership, private finance, community ownership, and public-private partnerships can significantly affect investment behavior, affordability, and public legitimacy.

I would conceptualize the UK's net-zero transition as a political-economic challenge involving four dimensions. First, ownership structures shape who controls low-carbon assets, who bears investment risk, and who captures the value created by transition. Second, governance systems determine whether decarbonization is coordinated across national government, local authorities, regulators, firms, and communities. Third, business models influence whether firms prioritise long-term public value, short-term profitability, affordability, or shareholder returns. Fourth, value logics shape how transition is justified: as a question of market efficiency, national competitiveness, social justice, resilience, or democratic renewal.

From this perspective, the key challenge is how to design institutions that make decarbonization socially acceptable, economically fair, and politically durable. This is especially important in sectors such as housing and transport, where transition policies affect everyday life and can generate resistance if costs are perceived as unfairly distributed. I would therefore be interested in examining alternative models of ownership and governance that can align decarbonization with affordability, accountability, and public legitimacy.

I conduct interdisciplinary research by combining industrial organization, computational economics, behavioral assumptions, policy analysis, and empirical modeling. My coauthored work on financial fraud detection also connects financial regulation, accounting, supervised machine learning, and investor protection. For this role, I would be keen to engage more deeply with political economy, economic history, sociology, sustainability transitions, and environmental governance. I see the ``new political economy of net zero'' as an agenda that should connect economic modeling with historically grounded and sociologically informed analysis.

My doctoral and coauthored research has involved extensive literature review and policy-oriented interpretation. My primary training is in economics and computational modeling, and I also have experience working with interpretive material through institutional analysis, policy comparison, and conceptual framework development. I am strongly motivated to further develop my qualitative research skills through this project, particularly through interviews with policymakers, industry stakeholders, and other actors involved in net-zero transitions. I would approach interviews as a way to understand how different actors interpret the constraints, conflicts, and possibilities of decarbonization, and to connect these insights to broader institutional mechanisms.

\section*{Communication, Collaboration, and Motivation}

I have presented my research at international conferences and workshops, including the International Conference on Challenges in Managing Smart Products and Services, the Annual Workshop on Economics with Heterogeneous Interacting Agents, and the German Network for New Economic Dynamics. These experiences have trained me to communicate complex models and technical findings to audiences with different disciplinary backgrounds.

I am confident in my ability to contribute to both academic publications and policy-relevant outputs. My writing aims to connect methodological rigour with practical relevance. In the context of this project, I would aim to communicate the political economy of net zero clearly to academic researchers, policymakers, and sectoral stakeholders, explaining how ownership, governance, and business models affect not only emissions outcomes, but also legitimacy, fairness, and institutional durability.

My academic training has taken place in an international and collaborative research environment across Germany and Italy. I have also worked on coauthored projects involving financial fraud detection and AI-supported research workflows. These experiences have taught me the importance of clear communication, transparent documentation, constructive feedback, and respect for different disciplinary perspectives. I would welcome the opportunity to work closely with Dr Donal Brown and the wider team at the Environmental Change Institute, contributing actively to collaborative publications, conceptual development, stakeholder engagement, and policy-facing outputs.

I am strongly motivated to apply for this role because it addresses one of the most important political-economic questions of the coming decades: how societies can decarbonise rapidly while building fairer, more legitimate, and more resilient economic institutions. The project's focus on ownership structures, governance systems, business models, and value logics is especially appealing because it goes beyond a narrow technical view of net zero. This position would be an excellent opportunity to extend my research from market and regulatory dynamics in digital sectors to the broader political economy of sustainability transitions and I am willing to contribute my skills and expertise.

\end{document}
